% Internal expected-result table. Values are target scenarios, not observations.
\documentclass[sigconf,anonymous]{aamas}
\usepackage{amsmath}
\usepackage{booktabs}
\usepackage{array}
\usepackage{makecell}
\usepackage{multirow}
\usepackage{xcolor}
\usepackage{colortbl}
\usepackage{url}
\usepackage{draftwatermark}
\input{aamas_metadata}

\vfuzz=1.5pt
\acmSubmissionID{INTERNAL EXPECTED VALUES -- NOT FOR SUBMISSION}
\submissionType{Research Paper Track}

\title[Expected Experiment Matrix]{Expected Experiment Matrix (Guide Values)}
\author{Anonymous Authors}
\affiliation{\institution{Internal research draft}}
\email{omitted}
\ccsdesc[500]{Computing methodologies~Multi-agent systems}
\ccsdesc[300]{Computing methodologies~Cooperation and coordination}
\keywords{multi-agent systems, role formation, peer judgment, online adaptation, experiments}

\newcommand{\blank}{\textemdash}
\newcommand{\unknown}{\textsc{unknown}}
\newcommand{\rare}{\textsc{rare}}
\newcommand{\metric}[1]{\makecell[c]{#1}}
\newcommand{\val}[2]{\makecell[c]{#1\\[-1pt]\scriptsize$\pm$#2}}
\newcommand{\svc}[2]{\makecell[c]{#1\\[-1pt]\scriptsize$\pm$#2}}
\newcommand{\pair}[2]{\makecell[c]{#1\\[-1pt]\scriptsize/#2}}
\SetWatermarkText{EXPECTED --- NOT OBSERVED}
\SetWatermarkColor{red!16}
\SetWatermarkAngle{35}
\SetWatermarkFontSize{28pt}

\begin{abstract}
This document is the numerical expected-result companion to the paper-facing experiment matrix. It preserves the original rows, columns, endpoints, and stream count so that every future observed cell has one target cell. Values are conservative scenario medians with expected stream variation, derived from the current mechanism conditions and public reference scales; they are not observations, fitted estimates, or evidence that a gate has passed.
\end{abstract}

\begin{document}
\maketitle

\section{Table set and correspondence}

The paper uses one compact overview and three result-bearing tables. The overview is a routing table. Table~\ref{tab:main-results} is the hero comparison for ArtifactRole; Table~\ref{tab:service-safety} reports service and safety; Table~\ref{tab:ablation} isolates mechanism boundaries. The expected table assumes three structural roots and three independent streams per arm ($n=9$); actual values may only replace these targets after the frozen cards, raw manifests and independent scoring pass.

\begin{table*}[t]
\centering
\small
\setlength{\tabcolsep}{3.5pt}
\renewcommand{\arraystretch}{1.12}
\begin{tabular}{p{0.08\textwidth}p{0.16\textwidth}p{0.23\textwidth}p{0.24\textwidth}p{0.20\textwidth}}
\toprule
RQ & Track / unit & Policy contrast & Primary result table & Primary estimand\\
\midrule
RQ1 & ArtifactRole / independent root & \rare{} vs raw acceptance, terminal-only, and contextual-trust-linear & Table~\ref{tab:main-results} & Common-cohort Brier/log loss and calibration\\
RQ2 & ArtifactRole / target assignment & Evidence plus delayed credit vs no evidence, delayed-only, and pooled & Tables~\ref{tab:main-results} and~\ref{tab:ablation} & Future quality--cost utility; change is diagnostic\\
RQ3 & PeerSelect / independent stream & Updater vs uniform, no-history, and history-only & Table~\ref{tab:service-safety} & Payoff/regret, update p95, backlog, state bytes\\
RQ4 & Both / mutation and drift streams & Ownership mutation, correction/replay, and leave-one-out & Tables~\ref{tab:service-safety} and~\ref{tab:ablation} & False attribution, coverage, forgetting/recovery\\
\bottomrule
\end{tabular}
\caption{Paper-facing experiment map. Each row points to the table that will carry the corresponding measured contrast. The table set is an overview of the pre-registered matrix; it is not evidence until the matching confirmation card is executed.}
\label{tab:overview}
\end{table*}

\paragraph{Expected-value rule.} Every number below is a target scenario, not an observed result. The displayed $\pm$ values are expected between-stream variation, not confidence intervals. The future experiment must retain all failures and UNKNOWN outcomes; it may disagree with this guide without being edited to match it.

\section{Primary result tables}

\subsection{ArtifactRole: information and assignment}

The first result table is the hero table. Its policy rows are the complete matched-information comparison boundary, while its columns are limited to the two primary scientific questions and complete cost. The row order is fixed before confirmation; the last row is the candidate mechanism, not a claim of superiority.

\begin{table*}[t]
\centering
\small
\setlength{\tabcolsep}{3pt}
\renewcommand{\arraystretch}{1.13}
\begin{tabular}{p{0.17\textwidth}*{6}{>{\centering\arraybackslash}p{0.105\textwidth}}}
\toprule
\multirow{2}{*}{Policy} & \multicolumn{2}{c}{RQ1: situated signal} & \multicolumn{2}{c}{RQ2: future assignment} & \multicolumn{1}{c}{Cost} & \multirow{2}{*}{Streams}\\
\cmidrule(lr){2-3}\cmidrule(lr){4-5}\cmidrule(lr){6-6}
 & \metric{Brier\\$\downarrow$} & \metric{Calibration\\$\downarrow$} & \metric{Quality--cost\\utility\\$\uparrow$} & \metric{\scriptsize Assignment\\[-1pt] change\\[-1pt] (diagnostic)} & \metric{Complete\\cost\\$\downarrow$} & \\
\midrule
Uniform & \val{0.250}{0.010} & \val{0.100}{0.015} & \val{0.400}{0.030} & 0.00 & \val{1.00}{0.04} & 9\\
No update & \val{0.238}{0.009} & \val{0.080}{0.013} & \val{0.430}{0.032} & 0.04 & \val{1.02}{0.04} & 9\\
Raw acceptance & \val{0.225}{0.010} & \val{0.073}{0.011} & \val{0.445}{0.034} & 0.10 & \val{1.10}{0.05} & 9\\
Terminal only & \val{0.220}{0.010} & \val{0.068}{0.010} & \val{0.452}{0.036} & 0.12 & \val{1.15}{0.05} & 9\\
Contextual-trust-linear & \val{0.211}{0.009} & \val{0.060}{0.009} & \val{0.466}{0.037} & 0.16 & \val{1.22}{0.06} & 9\\
Pooled controller & \val{0.207}{0.010} & \val{0.057}{0.010} & \val{0.458}{0.039} & 0.14 & \val{1.28}{0.06} & 9\\
\rare{} & \val{0.203}{0.009} & \val{0.055}{0.009} & \val{0.480}{0.038} & 0.20 & \val{1.34}{0.07} & 9\\
\bottomrule
\end{tabular}
\caption{Primary ArtifactRole result table. Every arm receives the same candidate menu, opportunity schedule, public-information boundary, selected-only rule, and complete-cost contract. Values are filled only after the independent root, later assignment, and outcome checks pass; intervals are stream-level intervals from the frozen card.}
\label{tab:main-results}
\end{table*}

\subsection{Online service and stability}

The second result table keeps operational claims separate from the hero utility comparison. PeerSelect supplies the controlled local-stream service measurement; ArtifactRole supplies the responsibility mutations. A method cannot trade away attribution safety for a lower update latency without showing that trade-off here.

\begin{table*}[t]
\centering
\small
\setlength{\tabcolsep}{2.6pt}
\renewcommand{\arraystretch}{1.12}
\begin{tabular}{p{0.16\textwidth}p{0.10\textwidth}*{6}{>{\centering\arraybackslash}p{0.095\textwidth}}}
\toprule
\multirow{2}{*}{Policy / intervention} & \multirow{2}{*}{Track} & \multicolumn{3}{c}{RQ3: service} & \multicolumn{3}{c}{RQ4: safety and stability}\\
\cmidrule(lr){3-5}\cmidrule(lr){6-8}
 & & \metric{Update\\p95\\$\downarrow$} & \metric{Backlog\\p95\\$\downarrow$} & \metric{State\\bytes\\(resource)} & \metric{False\\attribution\\$\downarrow$} & \metric{Coverage /\\UNKNOWN} & \metric{Forgetting /\\recovery}\\
\midrule
No update & PeerSelect & 0.00 & 0.0 & 8 & 0.18 & 0.55 & \pair{0.24}{11.0}\\
History only & PeerSelect & \svc{0.32}{0.08} & \svc{0.4}{0.2} & 20 & 0.14 & 0.61 & \pair{0.16}{8.0}\\
Contextual-trust-linear & PeerSelect & \svc{0.52}{0.12} & \svc{0.7}{0.3} & 28 & 0.13 & 0.64 & \pair{0.13}{7.0}\\
\rare{} & PeerSelect & \svc{0.68}{0.15} & \svc{1.0}{0.4} & 36 & 0.08 & 0.73 & \pair{0.10}{5.5}\\
Recipient-only mutation & ArtifactRole & \svc{0.30}{0.08} & \svc{0.1}{0.1} & 18 & 0.00 & 0.98 & \pair{0.00}{0.0}\\
Mixed ownership mutation & ArtifactRole & \svc{0.35}{0.10} & \svc{0.2}{0.1} & 22 & 0.02 & 0.92 & \pair{0.04}{1.0}\\
Delayed / out-of-order feedback & Both & \svc{0.74}{0.18} & \svc{1.4}{0.5} & 39 & 0.04 & 0.85 & \pair{0.12}{8.5}\\
\bottomrule
\end{tabular}
\caption{Service and stability table. Latency, backlog, and state size are measured under a fixed arrival and resource budget. Mutation rows test whether recipient-owned, mixed, delayed, or out-of-order events produce safe no-op or \unknown{} behavior rather than false producer credit.}
\label{tab:service-safety}
\end{table*}

\subsection{Mechanism and responsibility ablations}

The third table answers why any observed gain occurs. It holds the task roots, menus, models, and complete-cost accounting fixed while removing one information or update component at a time. These rows are not additional baselines; they are interventions on the candidate mechanism.

\begin{table*}[t]
\centering
\small
\setlength{\tabcolsep}{3pt}
\renewcommand{\arraystretch}{1.13}
\begin{tabular}{p{0.23\textwidth}p{0.15\textwidth}*{4}{>{\centering\arraybackslash}p{0.11\textwidth}}}
\toprule
\multirow{2}{*}{Intervention} & \multirow{2}{*}{\makecell[l]{Removed or changed\\component}} & \multicolumn{4}{c}{Primary endpoints}\\
\cmidrule(lr){3-6}
 & & \metric{Future\\utility\\$\uparrow$} & \metric{\scriptsize Assignment\\[-1pt] change\\[-1pt] (diagnostic)} & \metric{False\\attribution\\$\downarrow$} & \metric{Complete\\cost\\$\downarrow$}\\
\midrule
No evidence / no update & Public evidence and persistent update & \val{0.430}{0.032} & 0.04 & 0.18 & \val{1.02}{0.04}\\
Public evidence only & Persistent delayed update & \val{0.448}{0.034} & 0.12 & 0.11 & \val{1.18}{0.05}\\
Delayed update only & New public evidence & \val{0.455}{0.036} & 0.13 & 0.16 & \val{1.15}{0.05}\\
Evidence + delayed update (\rare{}) & None & \val{0.480}{0.038} & 0.20 & 0.08 & \val{1.34}{0.07}\\
\rare{} without ownership gate & Structural attribution gate & \val{0.461}{0.039} & 0.22 & 0.24 & \val{1.27}{0.06}\\
\rare{} without selected-only rule & Selected-only credit restriction & \val{0.472}{0.041} & 0.25 & 0.13 & \val{1.31}{0.07}\\
\rare{} without correction/replay & Version correction and replay checks & \val{0.468}{0.040} & 0.23 & 0.12 & \val{1.29}{0.07}\\
\bottomrule
\end{tabular}
\caption{Mechanism and responsibility ablations. The primary interpretation requires both a change in future assignment and a valid independent outcome. A utility gain without attribution safety, or a safety gain without changed assignment, is reported as a boundary result rather than role formation.}
\label{tab:ablation}
\end{table*}

\section{Population and result-entry contract}

The tables will be populated in a fixed order: first the independent-root and
same-information checks, then the primary ArtifactRole result, then the
PeerSelect service/stability table, and finally the ablation table. A table row
is reported as \unknown{} when its scorer, transport, or responsibility contract
cannot observe the required quantity. Failed attempts retain their measured cost;
they are not deleted to improve the mean. The final paper will add stream-level
intervals, per-root variation, and a short error analysis next to the relevant
table rather than hiding these quantities in a single aggregate number.

This blueprint is an internal pre-results artifact. It fixes presentation and
estimand alignment; it does not assert that any arm wins or that the benchmark
and baseline gates have passed.

\end{document}
